\begin{helper}
Fix a history cell \(H\), terminal coordinate set \(U\), nodup terminal
order \(o\), blue support label \(B\subseteq U\), and red label \(J\).  Let
\(\beta\) range over the complete terminal blue-trace labels supplied by
\(\noderef{FixedSetHistoryCellRedBranchTranscriptCells}\), and let
\(\tau=(P_\tau,A_\tau,Z_\tau)\) range over its finite transcript family.
Assume the terminal product law for \(H\), \(0\le p\le1/2\), the
fixed-\((B,J)\) realized-cell geometry exported by
\(\noderef{FixedSetHistoryCellRedBranchTranscriptCells}\), forward
compatibility with complete terminal assignments, and functional scan
disjointness.  Assume moreover explicit fixed-\((B,J)\) released-block data:
for every realized \((\beta,\tau)\) there is a finite block
\(\mathcal R_{\beta,\tau}\) of residual assignments \(A_{\rm res}\subseteq
U\).  For a residual assignment define its reconstructed full terminal
assignment by
\[
\operatorname{rec}_{\beta,\tau}(A_{\rm res})
  =(A_{\rm res}\setminus Z_\tau)\cup B\cup R_{\beta,J}.
\]
The released-block hypotheses say exactly the following.  Block assignments
are complete residual assignments on \(U\), membership in the block makes the
reconstructed full assignment compatible with \((\beta,\tau)\), reconstructed
compatibility puts the residual assignment back in the block, distinct
realized pairs have disjoint residual blocks, and
\(\mathcal R_{\beta,\tau}\) contains the whole residual cylinder obtained
from \(P_\tau,A_\tau\) after releasing the fixed supports
\(B\cup R_{\beta,J}\) and the selected absences \(Z_\tau\).  Thus this node
does not claim that the current RedBranch interface exports owner fibers; it
packages such block data when it is supplied as a separate fixed-\((B,J)\)
input.

Then the fixed terminal scan supplies explicit residual partition data.
There is a residual assignment mass \(\mu(A_{\rm res})\) on residual
terminal assignments \(A_{\rm res}\subseteq U\), released residual blocks
\(\mathcal R_{\beta,\tau}\), and released block masses
\(\rho_{\beta,\tau}\) such that:
\[
0\le \mu(A)\quad(A\subseteq U),\qquad
\sum_{A\subseteq U}\mu(A)\le1,
\]
for every realized pair,
\[
\mathcal R_{\beta,\tau}\subseteq\{A:A\subseteq U\},\qquad
\rho_{\beta,\tau}
  =\sum_{A\in\mathcal R_{\beta,\tau}}\mu(A),
\]
released blocks for distinct realized pairs are disjoint, and
\[
p^{|P_\tau|-u_R-u_B}(1-p)^{|A_\tau|-|Z_\tau|}
  \le \rho_{\beta,\tau}.
\]
\end{helper}
\begin{proof}
Use the released blocks supplied as explicit hypotheses for this fixed
\((B,J)\).  Those hypotheses are the deterministic scan-partition interface
needed here: each realized block consists of residual assignments on \(U\),
block membership implies compatibility of the reconstructed full assignment,
reconstructed compatibility implies block membership, and a residual
assignment cannot lie in two distinct realized blocks.  Thus the subset and
disjointness assertions for \(\mathcal R_{\beta,\tau}\) are part of the
input data of this helper, while the realized-cell geometry used below comes
from \(\noderef{FixedSetHistoryCellRedBranchTranscriptCells}\).

Define the mass on residual terminal assignments by the terminal Bernoulli
product weight
\[
\mu(A)=p^{|A|}(1-p)^{|U|-|A|}
\qquad(A\subseteq U),
\]
and put \(\mu(A)=0\) outside \(2^U\).  Since \(0\le p\le1/2\), both
\(p\) and \(1-p\) are nonnegative; hence \(\mu(A)\ge0\) for every
\(A\subseteq U\).  Applying
\(\noderef{FinsetPowersetCylinderMass}\) with empty present and absent
sets gives
\[
\sum_{A\subseteq U} p^{|A|}(1-p)^{|U|-|A|}
  =p^0(1-p)^0=1.
\]
Therefore \(\sum_{A\subseteq U}\mu(A)\le1\).  This is the complete terminal
product law written as a finite assignment mass; the conditional product law
hypothesis identifies the same weights with the conditional masses of the
complete terminal cylinders over \(H\).

For every formal pair set
\[
\rho_{\beta,\tau}
  =\sum_{A\in\mathcal R_{\beta,\tau}}\mu(A).
\]
This is exactly the required block-mass identity.  The block disjointness for
distinct realized pairs is one of the released-block hypotheses.  In the
intended deterministic scan application, that hypothesis is justified as
follows: if one residual assignment belonged to two realized blocks, then its
two reconstructed full assignments would be compatible with the two scan
certificates, and functional scan disjointness would force the two branch
labels and two transcripts to be equal.

It remains to prove the lower bound for a realized pair
\((\beta,\tau)\).  Let
\[
R=R_{\beta,J},\qquad
P^\circ=P_\tau\setminus(B\cup R),\qquad
A^\circ=A_\tau\setminus Z_\tau .
\]
The realized geometry from
\(\noderef{FixedSetHistoryCellRedBranchTranscriptCells}\) gives
\[
|R|=u_R,\qquad |B|=u_B,\qquad B\cup R\subseteq P_\tau,
\]
\[
Z_\tau\subseteq A_\tau,\qquad P_\tau\cap A_\tau=\emptyset,
\]
with \(B\) blue-only and \(R\) red-only.  Hence \(B\cap R=\emptyset\), so
\[
|P^\circ|=|P_\tau|-u_R-u_B.
\]
Also \(Z_\tau\subseteq A_\tau\), so
\[
|A^\circ|=|A_\tau|-|Z_\tau|.
\]
The disjointness \(P_\tau\cap A_\tau=\emptyset\) implies
\[
P^\circ\cap A^\circ=\emptyset .
\]

Consider the residual cylinder
\[
\mathcal C_{\beta,\tau}
  =\{A\subseteq U: P^\circ\subseteq A,\ A\cap A^\circ=\emptyset\}.
\]
The explicit released-block completeness hypothesis of this fixed
\((B,J)\) helper says precisely that
\(\mathcal C_{\beta,\tau}\subseteq\mathcal R_{\beta,\tau}\): after the
mandatory present supports \(B\cup R\) and the selected absences \(Z_\tau\)
are released, every residual assignment satisfying the remaining residual
present and absent constraints reconstructs to a full assignment accepted by
the same fixed \((B,J,\beta,\tau)\) scan block.

By \(\noderef{FinsetPowersetCylinderMass}\), applied to the disjoint
present set \(P^\circ\) and absent set \(A^\circ\), the \(\mu\)-mass of
\(\mathcal C_{\beta,\tau}\) is
\[
p^{|P^\circ|}(1-p)^{|A^\circ|}
 =
p^{|P_\tau|-u_R-u_B}(1-p)^{|A_\tau|-|Z_\tau|}.
\]
The same absent-factor accounting is the selected-absence split recorded in
\(\noderef{FixedSetHistoryCellSelectedAbsencePowSplit}\): the coordinates
of \(Z_\tau\) have been released once, so the residual absent exponent is
\(|A_\tau|-|Z_\tau|\).  Since
\(\mathcal C_{\beta,\tau}\subseteq\mathcal R_{\beta,\tau}\) and all
\(\mu\)-masses are nonnegative,
\[
p^{|P_\tau|-u_R-u_B}(1-p)^{|A_\tau|-|Z_\tau|}
  =\sum_{A\in\mathcal C_{\beta,\tau}}\mu(A)
  \le \sum_{A\in\mathcal R_{\beta,\tau}}\mu(A)
  = \rho_{\beta,\tau}.
\]
This proves the required residual domination and completes the construction
of the fixed-\((B,J)\) residual partition data.
\end{proof}
